%% ma: L12-B01-0018
%% lop: 12
%% chuong: 1
%% bai: 1
%% dang: 12.1.13
%% loai: DS
%% muc: [NB, TH, TH, TH]
%% dap_an: "ĐĐSĐ"
%% nguon: "Ảnh giáo viên gửi 11/10/2026 (Đề kiểm tra giữa HK1 Toán 12, THPT Việt Đức 2026-2027, mã đề 0103), Phần II Câu 1"
%% kien_thuc: "Bài 1 (cực trị, đơn điệu qua dấu đạo hàm), đạo hàm hàm lượng giác, phương trình lượng giác cơ bản (lớp 11)"
%% trang_thai: chinh-thuc
%% ghi_chu: "Đã duyệt 12/10/2026."
\begin{ex}
Xét hàm số $f(x)=x+2\cos x$ trên đoạn $[0;2\pi]$.
\choiceTF
{\True $f'(x)=1-2\sin x$.}
{\True Phương trình $f'(x)=0$ có hai nghiệm trong $(0;2\pi)$ là $x=\dfrac\pi6$ và $x=\dfrac{5\pi}6$.}
{Hàm số đạt cực đại tại điểm $x=\dfrac{5\pi}6$.}
{\True Hàm số nghịch biến trên khoảng $\left(\dfrac\pi6;\dfrac{5\pi}6\right)$.}
\loigiai{a) $f'(x)=1-2\sin x$. Đúng.

b) $f'(x)=0\Leftrightarrow\sin x=\dfrac12$; trong $(0;2\pi)$ có $x=\dfrac\pi6$ và $x=\dfrac{5\pi}6$. Đúng.

c) Trên $(0;2\pi)$: $f'(x)>0$ khi $x\in\left(0;\dfrac\pi6\right)\cup\left(\dfrac{5\pi}6;2\pi\right)$ và $f'(x)<0$ khi $x\in\left(\dfrac\pi6;\dfrac{5\pi}6\right)$. Do đó hàm số đạt cực đại tại $x=\dfrac\pi6$ và đạt cực tiểu tại $x=\dfrac{5\pi}6$. Khẳng định sai.

d) $f'(x)<0$ trên $\left(\dfrac\pi6;\dfrac{5\pi}6\right)$ nên hàm số nghịch biến trên khoảng đó. Đúng.}
\end{ex}
